\subsection{Additional Quantitative Results}
\label{sec:additional-quantitative-results}


% \documentclass{standalone}
% \usepackage{pgfplots}
% \usepgfplotslibrary{groupplots}
% \usetikzlibrary{patterns}

% \begin{document}

\begin{figure}[t]
\centering
\begin{tikzpicture}
    \begin{axis}[
        width=0.9*\linewidth,
        height=5cm,
        xlabel={Memory Horizon (s)},
        ylabel={Overall Level 2 AP (\%)},
        xmin=0, xmax=4.2,
        ymin=78, ymax=81.2,
        legend style={at={(0.98,0.5)}, anchor=east}, % Positioning the legend in the middle on the west side
        grid=major,
        ytick={78, 78.5, 79, 79.5, 80, 80.5, 81},
        xtick={0, 0.9, 1.8, 2.7, 3.6, 4.3}
    ]
    
    % Plot for 'SAFDNet 4f'
    \addplot[color=dark-color3, line width=1.5pt, dashed] coordinates {
        (0, 78.355)
        (0.9, 78.355)
        (1.8, 78.355)
        (2.4, 78.355)
        (2.7, 78.355)
        (3.6, 78.355)
        (4.2, 78.355)
    };
    \addlegendentry{\footnotesize{SAFDNet 4f~\cite{safdnet}}}
    
    % Plot for 'w/ Obj Mem'
    \addplot[color=dark-color1, mark=*, mark size=1.0pt, line width=1.5pt] coordinates {
        (0, 78.676)
        (0.9, 80.362)
        (1.8, 80.827)
        (2.4, 80.970)
        (2.7, 80.983)
        (3.6, 81.016)
        (4.2, 81.030)
    };
    \addlegendentry{\footnotesize{+ \ourmodel{}}}

    \end{axis}
\end{tikzpicture}% \end{document}
\caption{The effect of the memory horizon on the WOD validation set.
We keep the memory target time stride $s_m = 0.3s$ fixed,
while increasing the number of target timestamps $T_m \in \{0, 3, 6, 8, 9, 14\}$.
The base detector is SAFDNet 4f~\cite{safdnet}.}
\label{fig:history-curve}
    
\end{figure}
\subsubsection{Performance at different memory horizons}
In \cref{fig:history-curve},
we plot detection performance as a
function of the memory horizon $T_m s_m$ used at inference time. 
As described in \cref{sec:method-training}, the model is trained
with a variable number of target past timestamps $T_m \in \texttt{uniform}(\{6, 7, 8, 9, 10\})$,
with a variable stride $s_m \in \texttt{uniform}(\{0.2s, 0.3s, 0.4s\})$, meaning the
memory horizon during training varies from 1.2s to 4.0s.
%
To generate this plot, we run multiple evaluations on the WOD validation set,
each with a different number of target timestamps $T_m \in \{0, 3, 6, 8, 9, 14\}$,
keeping the memory target time stride $s_m = 0.3s$ fixed.
%
We find that the performance increases rapidly
when first introducing memory, and then increases more slowly and saturates
around 2.7s, due to the redundancy in most memory proposals.

\subsubsection{Performance across memory dropping rates}
% \documentclass{standalone}
% \usepackage{pgfplots}
% \usepgfplotslibrary{groupplots}
% \usetikzlibrary{patterns}

% \begin{document}

\begin{figure}[t]
\centering
    \begin{tikzpicture}
        \begin{axis}[
            width=0.9* \linewidth,
            height=5cm,
            xlabel={Memory Drop Rate},
            ylabel={Overall Level 2 AP (\%)},
            xmin=0, xmax=1,
            ymin=78, ymax=81.2,
            % legend pos=north east,
            legend style={at={(0.02,0.5)}, anchor=west},
            grid=major,
            ytick={78, 78.5, 79, 79.5, 80, 80.5, 81},
            xtick={0, 0.2, 0.4, 0.6, 0.8, 1}
        ]
        
        
        % Plot for '3D Detector'
        \addplot[color=dark-color3, dashed, line width=1.5pt, mark size=1pt] coordinates {
            (0, 78.4)
            (0.2, 78.4)
            (0.4, 78.4)
            (0.6, 78.4)
            (0.8, 78.4)
            (1, 78.4)
        };
        \addlegendentry{\footnotesize{SAFDNet 4f~\cite{safdnet}}}

        % Plot for 'w/ Obj Mem'
        \addplot[color=dark-color1, mark=*, line width=1.5pt, mark size=1pt] coordinates {
            (0, 80.97)
            (0.2, 80.87)
            (0.4, 80.83)
            (0.6, 80.55)
            (0.8, 79.84)
            (1, 78.68)
        };
        \addlegendentry{\footnotesize{+ \ourmodel{}}}
    
        \end{axis}
    \end{tikzpicture}

% \end{document}
\caption{The effect of randomly dropping memory target timestamps $\mathcal{T}_m$ at inference time on the WOD validation set.
The base detector is SAFDNet 4f~\cite{safdnet}.}
\label{fig:dropping-curve}
    
\end{figure}
To show the importance of the memory at inference time,
we investigate how the performance of \ourmodel{} changes
when we artificially drop memory proposals at inference time.
We randomly drop memory target past timestamps $\mathcal{T}_m$ with a given probability,
and \cref{fig:dropping-curve} plots model performance as a function of this dropping probability.
Our method is relatively robust up to a dropping probability of 0.6 because much of the information
provided by the memory proposals is redundant.
Removing all target timestamps $\mathcal{T}_m = \emptyset$ causes the performance to drop near the base detector, as expected.
The refinement transformer without any memory (dropping probability = 1.0) results
in \textit{slight} improvements over the base detector.
We hypothesize this is because our refinement transformer without memory
is the second stage of a two-stage detector,
which have been shown to improve performance~\cite{centerpoint}.
% \sergio{you can also point to Far3D. They use memory only for training and this brings them gains. I think it's because it adds hard examples for refinement }

\subsubsection{Forecasting Results}
\begin{table}[t]
    \centering
    \footnotesize
    \begin{tabular}{@{}lcccc@{}}
        \toprule
        Detector & Memory   & MR (\%) & ADE (m) & FDE (m) \\ 
        \midrule
        HEDNet 1f     & \xmark                   & 23.4          & 2.5     & 4.81    \\
        \rowcolor{tablecolor}
        HEDNet 1f     & \cmark                   & 18.1 & 0.89 & 2.05 \\ \midrule
        HEDNet 4f     & \xmark                   & 18.4          & 0.86    & 2.06    \\
        \rowcolor{tablecolor}
        HEDNet 4f     & \cmark                   & \textbf{17.2} & \textbf{0.72} & \textbf{1.75} \\
        \bottomrule
    \end{tabular}
    \caption{Measuring the effect of memory on trajectory forecasting performance on the Vehicle class in the WOD validation set.
    Metrics are computed on 5s forecasts at 0.5 Hz, at a recall threshold of 80\% with an IoU threshold of 0.5.
    Miss rate (MR) is computed at a distance threshold of 2\;m.}
    \label{tab:forecasting-metrics}
\end{table}

While our focus is task
of 3D object detection, \ourmodel{} also performs trajectory forecasting,
which is an important and well-studied sub-task of autonomous
driving systems~\cite{chai2019multipath,varadarajan2022multipath++,cui2019multimodal,phan2020covernet,cui2023gorela,ngiam2021scene}.
In \cref{tab:forecasting-metrics}, we investigate the effect of our memory mechanism on trajectory forecasting
performance on the WOD validation set by training
\ourmodel{} to enhance HEDNet 1f and HEDNet 4f~\cite{hednet} on the WOD validation set.
We also train a version where we remove the memory from the refinement transformer;
no memory proposals $\mathcal{P}^\mathrm{mem}$ (using $\mathcal{P}^{\mathrm{ref}(0)} = \mathcal{P}^\mathrm{det}$),
and no memory cross-attention, so that the model is still able to generate trajectory forecasts but cannot leverage memory inputs.
The metrics in \cref{tab:forecasting-metrics} are computed at a common detection recall point of 80\%
at an IoU threshold of 0.5.
Miss-rate (MR) measures the percentage of predicted trajectory endpoints that are more than 2m away from the ground truth trajectory endpoint.
Average displacement error (ADE) measures the average Euclidean distance between the predicted trajectory and the ground truth trajectory across all waypoints in the trajectory,
and final displacement error (FDE) measures the Euclidean distance between the predicted trajectory endpoint and the ground truth trajectory endpoint.
We find that memory improves trajectory forecasting performance for both HEDNet 1f and HEDNet 4f.
Note that, because HEDNet 4f uses multiple past LiDAR sweeps, it has better sensor evidence on the motion of vehicles.
This is why it results in better trajectory forecasting performance than HEDNet 1f;
the interpolated features $\mathbf{Q}^\mathrm{det}$ have more motion information which is useful to the refinement transformer.
Notably, HEDNet 1f with memory matches or outperforms HEDNet 4f without memory, showing the strength of our memory mechanism in capturing object motion information.


\subsubsection{Computational Overhead}

We report total FLOPs and peak memory measurements per forward pass,
taking the maximum across 10 random sequences on WOD in \cref{tab:flops}.
\begin{table}
    \begin{center}
    \footnotesize
    \begin{tabular}{lcc}
    \toprule
    \textbf{Model} & \textbf{GigaFLOPs} & \textbf{GPU Mem (GB)} \\
    \midrule
    SAFDNet 4f        & 880 & 2.5 \\
    \ourmodel{} Module        & 50  & 2.8 \\
    \bottomrule
    \end{tabular}
    \end{center}
    \caption{FLOPs and peak memory measurements per forward pass for SAFDNet 4f and \ourmodel{}.}
    \label{tab:flops}
\end{table}
We use FLOPs because it is more hardware and implementation agnostic.
There is minimal computational overhead from \ourmodel{} relative to the base detector.
We also note that in most traditional self-driving systems, the 3D detector would be followed by a forecasting module,
which \ourmodel{} can replace.

\subsubsection{Performance on Stationary Objects}

\Cref{fig:stationary-objects} shows the performance of \ourmodel{} on stationary objects on WOD ($\le 0.5$ m/s).
\begin{table}
\begin{center}
\footnotesize
\begin{tabular}{lccc}
\toprule
 & \multicolumn{3}{c}{\textbf{Stationary}} \\
 & \textbf{Veh. AP} & \textbf{Ped. AP} & \textbf{Cyc. AP} \\
\midrule
SAFDNet 4f (Base)          & 73.8 & 54.1 & 34.3 \\
\ourmodel{} No Forecasting & \textbf{74.9} & \textbf{56.2} & 35.9 \\
\ourmodel{}                & 74.8 & 56.1 & \textbf{36.9} \\
\bottomrule
\end{tabular}
\end{center}
\caption{Performance on stationary objects on WOD ($\le 0.5$ m/s).}
\label{fig:stationary-objects}
\end{table}
As expected, forecasting provides minimal benefit for detecting stationary objects, and both versions of \ourmodel{} provide improvements over the base detector.